\documentclass[11pt]{article}
\usepackage[a4paper,margin=28mm]{geometry}
\usepackage{amsmath,amssymb}
\usepackage{hyperref}
\setlength{\emergencystretch}{2em}
\title{An infinite prime chain inside a localization of a PI ring}
\author{Ziang Ni (\texttt{ziangni-sys})}
\date{10 October 2026}
\begin{document}
\maketitle
\noindent\textbf{Disproof of conjecture 00000009976.}
The conjecture asserts finiteness of prime localization chains,
with a bound polynomial in the polynomial identity degree. A
commutative polynomial ring provides an actual infinite chain in a
single localization, with fixed identity degree.

Let
\[
 R=\mathbb Q[X_0,X_1,X_2,\ldots].
\]
This ring is commutative, so it satisfies the nonzero polynomial
identity $UV-VU=0$ of degree two in the free associative algebra.
Thus it is a PI ring. If ``PI degree'' denotes the conventional PI
degree of a commutative domain, that degree is one instead; either
convention gives a fixed finite degree here.

For each $n\geq0$, let $\kappa_n:R\to R$ be the substitution
\[
 \kappa_n(X_i)=\begin{cases}0&i<n,\\X_i&i\geq n,\end{cases}
 \qquad P_n=\ker\kappa_n.
\]
Because the target ring is an integral domain, $P_n$ is prime.
Equivalently $P_n=(X_0,\ldots,X_{n-1})$, and its quotient is the
polynomial ring in the remaining variables. If $m\leq n$ then
$\kappa_n\kappa_m=\kappa_n$, hence $P_m\subseteq P_n$.
Moreover
\[
 X_n\in P_{n+1},\qquad X_n\notin P_n.
\]
Thus $P_0\subsetneq P_1\subsetneq P_2\subsetneq\cdots$ is an
infinite strict chain of prime ideals.

Let $\varepsilon:R\to\mathbb Q$ evaluate every $X_i$ at zero, and
let $M=\ker\varepsilon$. It is prime (indeed maximal, since
$\varepsilon$ is surjective), and $\varepsilon\kappa_n=\varepsilon$
shows that every $P_n\subseteq M$.

Now use the actual localization $R_M=(R\setminus M)^{-1}R$.
Prime localization correspondence is the order isomorphism
\[
 \{\text{prime ideals of }R_M\}\ \cong\
 \{\text{prime ideals of }R\text{ contained in }M\}.
\]
Its inverse extends an ideal to $R_M$, and contraction recovers the
original prime ideal. Consequently $Q_n=P_nR_M$ is prime and
\[
 Q_0\subsetneq Q_1\subsetneq Q_2\subsetneq\cdots.
\]
For completeness, strictness cannot disappear: if $Q_n=Q_{n+1}$,
contracting both sides would give $P_n=P_{n+1}$, contradicting the
variable witness above. This is an infinite prime chain inside a
single localization of a ring with fixed PI degree. In particular
there are finite initial chains of every length, so no finite
polynomial bound in that degree can hold.

The asserted chain finiteness is therefore false. We need not
dispute the separate spectral-space claim or assign a meaning to
the source's unspecified comparison topology: its explicit
finiteness clause already fails. A Noetherian or finite-dimensional
restriction would be additional to the bilingual source, which
imposes neither.

\paragraph{Lean coverage.}
The project uses the actual ring of multivariate polynomials in
countably many variables over $\mathbb Q$. It constructs substitution
homomorphisms, their prime kernels, containment in the same prime
$M$, and strictness from actual variables. Crucially it then uses
the standard library's prime order isomorphism for the actual ring
\texttt{Localization.AtPrime M}, producing a formal infinite strict
chain of localized prime ideals. It also proves the fixed degree-two
commutator identity. Lean 4.19.0 and Mathlib are publicly pinned;
printed theorem audits use only standard logical axioms.
\end{document}
